\subsection{Demonstrated Scientific Gains and Their Attribution}
\label{sec:comparative-gains}

The primary comparisons establish several kinds of gain: more informative experiment selection, improved candidate performance, better search procedures, and more effective access to scientific tools. Their attribution depends on the comparator. A comparison with an initial formulation establishes improvement during a campaign; a comparison with another selection policy addresses how experiments are chosen; a comparison with the same underlying language model addresses the value of the added system. Table~\ref{tab:comparative-gains} separates these interpretations.

\begin{table}[H]
\caption{Comparative evidence for scientific gains in selected primary studies. Metrics retain their study-specific meanings and are not pooled across domains.}
\label{tab:comparative-gains}
\scriptsize
\begin{tabularx}{\textwidth}{@{}p{0.14\textwidth}>{\raggedright\arraybackslash}X>{\raggedright\arraybackslash}X>{\raggedright\arraybackslash}X@{}}
\toprule
Study & Comparator and endpoint & Reported result & Scope of attribution \\
\midrule
BEAR~\cite{asd0504} & Adaptive physical design search versus a benchmark derived from 1,800 grid samples; measured toughness & A batch campaign matched benchmark toughness after 32 experiments (12 h) and exceeded it after 64 (24 h) & Campaign-level selection efficiency within one structural family; benchmark construction is an additional cost \\
Mobile robotic chemist~\cite{burger2020} & Final versus initial formulation; hydrogen-evolution rate & 688 experiments in eight days; approximately sixfold improvement & Joint outcome of automated experimentation and adaptive search; no sixfold advantage over a matched human or random-search control \\
FunSearch~\cite{asd0512} & Evolved versus first-fit and best-fit heuristics; excess bins relative to a lower bound & On OR4: 2.47\% excess bins versus 5.23\% and 4.94\%, respectively & Better tested packing heuristics; the percentages do not measure search cost or the number of bins saved directly \\
Bio\-Discovery\-Agent~\cite{asd0020} & Claude 3.5 Sonnet without tools versus the best tabulated baseline per dataset; five-round hit ratio & CAR-T: 0.130 versus 0.069; Carnevale: 0.042 versus 0.047; improvement on five of six datasets & Retrospective retrieval of perturbation outcomes; comparison includes pretrained biological knowledge \\
ChemCrow~\cite{bran2024} & Tool-assisted system versus GPT-4 without tools; expert assessment of chemistry tasks & Expert assessments favored tool-assisted responses & Supports the system-level benefit of tool access and orchestration; does not isolate each component's contribution \\
\bottomrule
\end{tabularx}
\end{table}

Experiment-selection efficiency has evidence predating LLM-based systems. The Robot Scientist compared its cost-aware selection strategy with cheaper-first and random alternatives, evaluating the accuracy of remaining hypotheses within the supplied biological model~\cite{king2004}. BEAR provides a physical engineering comparison: adaptive campaigns identified high-toughness structures with far fewer trials than were used to construct the grid benchmark~\cite{asd0504}. The study expended more than 2,500 experiments in establishing and examining this comparison. The 32-trial result therefore characterizes the demonstrated search campaign, rather than the total cost of creating the platform and its benchmark. Together, these studies show how selecting informative experiments can improve research efficiency independently of language-model reasoning.

Improved measured performance supports a different conclusion from superiority of the selection method. The mobile robotic chemist's approximately sixfold increase in hydrogen evolution establishes successful optimization within its formulation space~\cite{burger2020}. Without a matched alternative campaign, that ratio cannot apportion the improvement among throughput, Bayesian selection, and the researcher-designed ingredient space. Similarly, Granda and colleagues combined automated reactivity screening with manual follow-up, isolation, and characterization of selected transformations~\cite{asd0734}. The scientific result depends on both the screening process and the subsequent work that establishes the chemical identity and significance of its outputs. These examples reveal a contribution distributed across investigation stages, rather than a single algorithmic gain measured by the final artifact.

FunSearch demonstrates an additional route to contribution: generating a reusable procedure that yields improved solutions~\cite{asd0512}. Its bin-packing comparison tests the resulting heuristic against established methods on separate benchmark instances. Its mathematical results also include new constructions, with further human analysis of discovered structure enabling additional improvements. The LLM supplies program variations; evolutionary selection and executable evaluation determine which variations survive. This combination expands the search over candidate procedures, while the reported heuristic advantage concerns the completed system's output. It does not imply that an unaided LLM solves the mathematical problem or that greater workflow breadth caused the improvement.

BioDiscoveryAgent addresses whether pretrained biological knowledge can improve sequential selection relative to data-driven acquisition baselines~\cite{asd0020}. Its five-round evaluation selects batches of 128 genes, reduced to 32 for the smaller dataset, and retrieves their measured outcomes from past screens. Claude 3.5 Sonnet without tools exceeds the strongest tabulated baseline on five of six datasets, with substantial variation between biological contexts (Table~\ref{tab:comparative-gains}). The early-round advantage is consistent with informative prior knowledge mitigating cold start. Because the LLM and the acquisition baselines have different prior information, the result does not isolate feedback-based reasoning from pretrained knowledge. Component comparisons further show that added tools improve the weaker Claude 3 Haiku model more consistently than Claude 3.5 Sonnet. Incorporating prior knowledge and adding orchestration are therefore distinguishable contributions whose benefits depend on the model and task.

ChemCrow's comparison with tool-free GPT-4 more directly tests an extension to the same base model~\cite{bran2024}. Expert preference supports the added system's utility for the tested chemistry tasks, while leaving retrieval, specialized algorithms, and orchestration jointly responsible. Across the cases, demonstrable gains attach to particular decisions, outputs, and resource conditions. A stronger claim about a decision paradigm requires a comparator that preserves the available information, tools, experimental budget, and human contribution relevant to that claim. The current evidence supports complementary contributions from symbolic hypothesis testing, optimization, planning, learned policies, and LLM-based integration; it does not establish a universal advantage for one architecture.
